\section{Avaliação da não-linearidade através do princípio da superposição}

\begin{frame}{O Princípio da Superposição}
  Para determinar a natureza do sistema, aplicamos o \textbf{Princípio da Superposição}, que define um sistema linear se:
  \begin{equation}
    y(u_1 + u_2) = y(u_1) + y(u_2)
  \end{equation}

  A não-linearidade é devido a:
  \begin{itemize}
    \item \textbf{Cinética de Arrhenius:} O termo $k = k_0 \exp(-E/RT)$ introduz dependência exponencial da temperatura.
    \item \textbf{Balanços de Energia:} Interações entre fluxos e temperaturas nos denominadores e produtos de estados.
  \end{itemize}
  \vspace{0.5cm}

\end{frame}

\begin{frame}{Metodologia de Avaliação}
  O teste foi realizado no código através de três cenários de perturbação:
  \begin{itemize}
    \item \textbf{Teste 1 ($\Delta y_1$):} Degrau de $+10\%$ no calor do reator 1 ($Q_1$).
    \item \textbf{Teste 2 ($\Delta y_2$):} Degrau de $+10\%$ no fluxo de alimentação ($F_{f1}$).
    \item \textbf{Teste Combinado ($\Delta y_{total}$):} Aplicação simultânea de $Q_1$ e $F_{f1}$ ($+10\%$).
  \end{itemize}
  \vspace{0.3cm}
  O \textbf{Erro de Superposição} mede o desvio do comportamento linear:
  \begin{equation}
    \text{Erro}(t) = \Delta y_{total}(t) - (\Delta y_1(t) + \Delta y_2(t))
  \end{equation}
\end{frame}

\begin{frame}{Resultado da Avaliação de Não Linearidade}
  \begin{figure}
    \centering
    \includegraphics[width=0.5\linewidth]{figuras/Grafico.png}
    \caption{Colocar legenda aqui}
  \end{figure}
\end{frame}
